\chapter{系统需求分析}
\label{chap:requirement}

需求分析是连接课题目标与系统设计的中间环节。本章从可行性、用户角色、功能需求、非功能需求和业务流程五个方面展开，明确系统“为谁做、做什么、做到什么程度”，为第~\ref{chap:design} 章的设计提供依据。需求来源包括任务书、开题报告、校园常见问题以及对学生和管理人员的访谈记录。

\section{可行性分析}
\label{sec:feasibility}

\subsection{技术可行性}

系统所需的技术组件均有成熟方案。前端采用 Vue 3 与 Element Plus，能够快速构建表单、列表、对话框和聊天界面；后端采用 Python 与 FastAPI，异步框架适合处理文档解析、向量检索和大模型调用等 I/O 密集型任务；MySQL 适合保存结构化业务数据；Chroma 提供轻量的向量存储和相似度检索接口；BGE 中文嵌入模型和交叉编码器可以在 CPU 上运行，也可以按需替换为其他模型；大模型通过 HTTP 接口调用，不需要本地显卡。上述组件均为开源或公开可用的服务，文档齐全，社区活跃，出现问题时容易找到解决方案。

从开发能力看，本人在本科阶段学习了软件工程、数据库、Web 开发、数据结构与算法等课程，能够完成前后端分离项目的设计与开发；了解 Python、JavaScript、SQL 和 RESTful 接口的基本用法；通过文献阅读和实验，掌握了文档解析、文本分块、向量化、检索和提示词构造的基本流程。课题不要求训练或微调大模型，技术风险主要集中在文档解析质量和检索效果调优上，这两部分都可以通过分阶段迭代解决。

\subsection{经济可行性}

系统开发和运行的成本较低。开发阶段使用个人计算机和开源软件，不需要购买服务器和商业数据库；嵌入模型和重排序模型可以在本地 CPU 上运行，不需要 GPU；大模型按调用量计费，测试阶段的使用量有限。部署阶段可以运行在学校已有的虚拟机上，也可以使用低成本的云主机；向量数据库和 MySQL 均可单机部署，硬件要求不高。与采购商业问答系统相比，本方案的一次性投入和长期维护成本都更容易控制。

\subsection{操作可行性}

系统采用 B/S 架构，用户通过浏览器即可访问，不需要安装客户端。学生端界面以对话为主，输入问题后直接查看答案和引用来源，操作路径短；管理端采用常见的后台布局，资料上传、分类、知识库构建和统计查看都以表单和列表呈现，符合管理人员的使用习惯。系统提供明确的错误提示和空状态引导，例如检索无结果时提示用户换一种问法，文档解析失败时显示失败原因和重试按钮，降低使用门槛。

\subsection{进度可行性}

按照任务书的进度安排，2026 年 9 月至 10 月完成调研、开题和需求分析，2026 年 10 月至 11 月完成核心模块开发，2026 年 12 月至 2027 年 2 月完成测试、论文和修改，2027 年 3 月完成查重、答辩和归档。本课题的功能范围明确，技术路线成熟，关键模块可以并行开发：文档处理和知识库构建、检索与重排序、问答与会话、管理统计相对独立，便于分阶段联调。只要控制好知识库范围和测试集规模，进度风险可控。

\section{用户角色与权限需求}
\label{sec:roles}

系统面向三类用户。学生是主要使用者，可以登录、提问、查看引用、继续追问、查询历史和提交满意度反馈，不能上传或删除知识资料。教师可以查看与自己课程或业务相关的问答统计，协助维护常见问题，但不具备用户管理和系统配置权限。管理员负责资料管理、知识库构建、用户与角色管理、统计分析和参数配置，是权限最高的角色。为避免权限过大，系统可以设置一名超级管理员，只用于初始化和应急处理，日常管理由普通管理员完成。

角色与功能的对应关系如表~\ref{tab:role-matrix} 所示。“查看”表示只能读取，“维护”表示可以新增、修改和删除，“管理”表示可以分配权限或修改系统级配置。权限控制采用“角色—权限—接口”三级映射：角色绑定权限集合，权限对应后端接口，接口在处理请求前校验当前用户是否具备相应权限。前端根据权限隐藏菜单和按钮，但后端校验才是最终防线。

\begin{table}[H]
  \centering
  \caption{角色与功能权限矩阵}
  \label{tab:role-matrix}
  \zihao{5}
  \renewcommand{\arraystretch}{1.25}
  \begin{tabular}{lcccc}
    \toprule
    功能模块 & 学生 & 教师 & 管理员 & 超级管理员 \\
    \midrule
    登录与个人信息 & 使用 & 使用 & 使用 & 使用 \\
    智能问答与追问 & 使用 & 使用 & 使用 & 使用 \\
    引用来源查看 & 查看 & 查看 & 查看 & 查看 \\
    问答历史与反馈 & 维护 & 维护 & 查看 & 查看 \\
    资料上传与分类 & 无 & 无 & 维护 & 维护 \\
    知识库构建与重建 & 无 & 无 & 维护 & 维护 \\
    高频问题与统计 & 无 & 查看 & 查看 & 查看 \\
    用户与角色管理 & 无 & 无 & 管理 & 管理 \\
    系统参数与密钥配置 & 无 & 无 & 维护 & 管理 \\
    \bottomrule
  \end{tabular}
\end{table}

除功能权限外，知识资料还需要内容级权限。校园公开制度、办事指南可以对所有登录用户开放；涉及教师评价、学生资助名单等敏感内容的文档，应限制在特定角色范围内。检索时先按用户角色过滤候选片段，再计算相似度，避免低权限用户通过问答间接获取受限信息。文档的权限标签由管理员在 upload 时设置，默认继承所属分类的权限，特殊文档可以单独覆盖。

\section{功能需求分析}
\label{sec:functional}

根据任务书和开题报告，系统需要实现用户管理、资料管理、知识库构建、智能问答、引用溯源、会话反馈和统计分析等功能。下面按“输入—处理—输出—异常”的方式逐项说明，功能编号与后续设计和测试用例保持一致。

（1）用户登录与权限管理（FR-01）。输入为账号、密码和验证码；处理时校验账号状态、密码哈希和验证码，签发 JWT，并根据角色加载权限；输出为令牌、用户信息和菜单权限；异常包括账号不存在、密码错误、账号被禁用、令牌过期，系统应返回明确的错误码并记录登录日志。

（2）校园资料上传与分类（FR-02）。输入为 PDF、Word 或纯文本文件，以及标题、分类、权限标签和备注；处理时校验文件类型和大小，保存原文件，创建文档记录，触发解析任务；输出为文档编号和解析状态；异常包括格式不支持、文件损坏、大小超限、重复上传，系统应拒绝或提示，并允许用户重试。

（3）文档解析、分块与向量化（FR-03）。输入为待处理的文档；处理时提取文本，清洗页眉页脚和多余空白，按标题层级和段落结构分块，保留页码和标题路径，调用嵌入模型生成向量，写入向量数据库；输出为知识片段和索引；异常包括解析失败、编码错误、模型服务不可用，系统应记录失败原因，支持断点重试和幂等重建。

（4）自然语言问答（FR-04）。输入为用户问题、会话编号和可选的分类过滤条件；处理时进行问题改写、双路召回、RRF 融合、交叉编码器重排，组装提示词并调用大模型；输出为答案、引用编号、候选证据和耗时；异常包括检索为空、模型超时、模型返回格式错误，系统应降级为拒答或提示稍后重试，不能返回无依据的内容。

（5）引用溯源与来源展示（FR-05）。输入为生成结果中的引用编号；处理时把编号映射回知识片段，校验片段是否属于本次候选集合，截取原文摘要；输出为可展开的来源卡片，包含文档标题、章节路径、页码和原文片段；异常包括编号越界、片段被删除、权限变化，系统应隐藏无效引用并提示用户。

（6）问答历史、继续提问与满意度反馈（FR-06）。输入为用户编号、会话编号和反馈值；处理时保存会话、消息、引用和反馈，支持按时间、关键词和分类查询；输出为历史列表、会话详情和反馈统计；异常包括会话不存在、无权查看、重复提交，系统应做幂等处理并返回统一错误信息。

（7）管理统计与系统配置（FR-07）。输入为时间范围、分类和统计维度；处理时聚合高频问题、未解决问题、知识库使用量、满意度和平均响应时延；输出为图表和明细列表；配置项包括检索条数、融合常数、重排阈值、模型名称、超时时间和开关；异常包括时间范围非法、配置值越界，系统应拒绝保存并给出取值范围。

上述功能的优先级和验收标准如表~\ref{tab:function-list} 所示。优先级“高”表示首版必须完成，“中”表示在核心链路稳定后完成，“低”表示可在后续版本迭代。

\begin{table}[H]
  \centering
  \caption{功能需求清单}
  \label{tab:function-list}
  \zihao{5}
  \renewcommand{\arraystretch}{1.25}
  \begin{tabular}{p{1.3cm}p{3.2cm}p{6.8cm}p{1.4cm}}
    \toprule
    编号 & 功能名称 & 主要验收标准 & 优先级 \\
    \midrule
    FR-01 & 登录与权限管理 & 正确签发和校验令牌；越权接口返回 403；登录失败有提示 & 高 \\
    FR-02 & 资料上传与分类 & 支持 PDF、DOCX、TXT；分类和权限标签可编辑；重复文件有提示 & 高 \\
    FR-03 & 解析、分块与向量化 & 文本提取完整；分块保留标题和页码；失败可重试 & 高 \\
    FR-04 & 自然语言问答 & 8 秒内返回；回答带引用；证据不足时拒答 & 高 \\
    FR-05 & 引用溯源与展示 & 引用可定位到文档、章节和页码；无效引用不展示 & 高 \\
    FR-06 & 历史、追问与反馈 & 历史可查询；多轮追问能理解指代；反馈可统计 & 中 \\
    FR-07 & 统计与系统配置 & 高频问题、未解决问题、满意度可查看；配置有范围校验 & 中 \\
    \bottomrule
  \end{tabular}
\end{table}

\section{非功能需求分析}
\label{sec:nonfunctional}

非功能需求决定系统是否“好用”和“可维护”。结合校园问答的使用场景，本文提出以下要求。

（1）性能。单次问答的平均响应时间不超过 8 秒，其中检索与重排不超过 2 秒，大模型生成不超过 6 秒；文档解析和向量化采用异步任务，不阻塞用户操作；管理端统计查询在 3 秒内返回。对于超过 8 秒的请求，前端显示加载状态和取消按钮，后端记录超时原因。

（2）并发与可用性。系统面向校内用户，按 50 个并发问答请求设计；采用无状态后端和连接池，支持水平扩展；关键服务出现故障时，问答接口降级为“仅返回检索片段”，保证用户仍能查看资料；系统可用性目标为 99\%。

（3）可靠性。文档处理、向量化和模型调用都可能失败，系统需要记录任务状态、失败原因和重试次数，支持手动重试和定时补偿；问答记录、反馈和引用在事务中保存，不能因为一次模型调用失败而丢失；向量库和 MySQL 的数据需要定期备份。

（4）安全性。所有接口除登录外都需要认证；密码加盐哈希存储；大模型密钥和数据库密码通过环境变量注入；上传文件做类型、大小和内容校验；检索结果按权限过滤；日志脱敏；对登录、上传、问答等接口做频率限制，防止暴力破解和资源滥用。

（5）可维护性。前后端分离，接口使用统一的响应结构和错误码；业务逻辑按模块划分，模型服务、向量库和文档解析器都通过接口抽象，便于替换；关键操作记录结构化日志，日志包含请求编号、用户编号、耗时和错误码，方便排查。

（6）可扩展性。知识库支持新增分类和文档类型；检索参数、模型名称和提示词模板可以在管理端配置；向量库、嵌入模型和大模型服务可以替换；后续可以接入知识图谱、语音输入和多模态文档，而不需要重写问答主流程。

（7）易用性。界面以对话为核心，问题输入框支持多行和快捷发送；答案中的引用编号可点击展开；空状态、加载状态和错误状态都有明确提示；管理端表格支持筛选、分页和批量操作。

\section{用例与业务流程}
\label{sec:usecase}

系统的用例关系如图~\ref{fig:usecase} 所示。学生的主要用例包括登录、提问、查看引用、继续追问、查询历史和提交反馈；教师可以查看统计和协助维护常见问题；管理员在此基础上增加资料管理、知识库构建、用户角色管理和系统配置。所有用户都需要先登录，问答和引用查看依赖于知识库中存在可检索的资料。

\begin{figure}[H]
  \centering
  \begin{tikzpicture}[
    uc/.style={draw, ellipse, minimum width=2.9cm, minimum height=0.7cm, align=center, fill=gray!8, font=\zihao{5}},
    actor/.style={font=\zihao{5}},
    line/.style={thick}
  ]
    % 系统边界
    \draw[rounded corners=4pt, thick] (0.8,-0.9) rectangle (9.4,5.4);
    \node[font=\zihao{5}\heiti] at (5.1,5.15) {校园智能问答系统};

    % 学生
    \node[actor] (stu) at (-1.2,3.0) {学生};
    \draw[line] (-1.2,2.65) circle (0.18);
    \draw[line] (-1.2,2.47) -- (-1.2,1.65);
    \draw[line] (-1.55,2.2) -- (-0.85,2.2);
    \draw[line] (-1.2,1.65) -- (-1.55,1.0);
    \draw[line] (-1.2,1.65) -- (-0.85,1.0);

    % 管理员
    \node[actor] (adm) at (11.0,3.0) {管理员};
    \draw[line] (11.0,2.65) circle (0.18);
    \draw[line] (11.0,2.47) -- (11.0,1.65);
    \draw[line] (10.65,2.2) -- (11.35,2.2);
    \draw[line] (11.0,1.65) -- (10.65,1.0);
    \draw[line] (11.0,1.65) -- (11.35,1.0);

    % 左侧：学生用例
    \node[uc] (login) at (3.0,4.6) {登录与权限};
    \node[uc] (ask) at (3.0,3.5) {自然语言问答};
    \node[uc] (cite) at (3.0,2.4) {引用来源查看};
    \node[uc] (history) at (3.0,1.3) {历史与反馈};

    % 右侧：管理员用例
    \node[uc] (upload) at (7.4,4.6) {资料上传与分类};
    \node[uc] (build) at (7.4,3.5) {知识库构建};
    \node[uc] (stat) at (7.4,2.4) {统计分析};
    \node[uc] (config) at (7.4,1.3) {系统配置};
    \node[uc] (user) at (7.4,0.2) {用户角色管理};

    % 学生连线
    \draw[line] (-0.9,3.0) -- (login.west);
    \draw[line] (-0.9,3.0) -- (ask.west);
    \draw[line] (-0.9,3.0) -- (cite.west);
    \draw[line] (-0.9,3.0) -- (history.west);

    % 管理员连线
    \draw[line] (10.7,3.0) -- (upload.east);
    \draw[line] (10.7,3.0) -- (build.east);
    \draw[line] (10.7,3.0) -- (stat.east);
    \draw[line] (10.7,3.0) -- (config.east);
    \draw[line] (10.7,3.0) -- (user.east);
  \end{tikzpicture}
  \caption{系统用例图}
  \label{fig:usecase}
\end{figure}

问答业务流程如图~\ref{fig:qa-flow} 所示。用户提交问题后，系统先结合最近会话做问题改写，再分别执行 BM25 召回和向量召回，用 RRF 合并候选并重排。若最高相关性低于阈值或候选为空，系统不调用大模型，直接返回“暂时找不到依据”的提示，并给出可能的分类建议。若证据充足，系统把问题、证据和输出约束组装成提示词，调用大模型生成答案。生成后再做引用校验：引用编号必须落在候选证据范围内，答案中的关键结论必须能对应到至少一条证据。校验通过则返回答案和来源卡片；校验不通过则降级为摘录式回答，或提示用户补充信息。

\begin{figure}[H]
  \centering
  \begin{tikzpicture}[
    node distance=0.55cm,
    flowstep/.style={draw, rounded corners=3pt, minimum width=6.4cm, minimum height=0.75cm, align=center, fill=gray!8, font=\zihao{5}},
    dec/.style={draw, diamond, aspect=2.4, minimum width=5.6cm, minimum height=1.0cm, align=center, fill=gray!16, font=\zihao{5}, inner sep=1pt},
    arr/.style={-{Latex[length=2mm]}, thick}
  ]
    \node[flowstep] (q) {用户输入问题};
    \node[flowstep, below=of q] (rw) {结合会话历史改写问题};
    \node[flowstep, below=of rw] (ret) {BM25 与向量双路召回};
    \node[flowstep, below=of ret] (fu) {RRF 融合与交叉编码器重排};
    \node[dec, below=of fu] (enough) {证据是否充足};
    \node[flowstep, below=1.0cm of enough] (gen) {受约束生成答案与引用};
    \node[dec, below=of gen] (valid) {引用是否有效};
    \node[flowstep, below=1.0cm of valid] (ok) {返回答案与来源卡片};
    \node[flowstep, right=1.0cm of enough, minimum width=3.2cm] (refuse) {拒答并给出建议};
    \node[flowstep, right=1.0cm of valid, minimum width=3.2cm] (fallback) {降级为摘录回答};

    \draw[arr] (q) -- (rw);
    \draw[arr] (rw) -- (ret);
    \draw[arr] (ret) -- (fu);
    \draw[arr] (fu) -- (enough);
    \draw[arr] (enough) -- node[left, font=\zihao{6}] {是} (gen);
    \draw[arr] (enough) -- node[above, font=\zihao{6}] {否} (refuse);
    \draw[arr] (gen) -- (valid);
    \draw[arr] (valid) -- node[left, font=\zihao{6}] {是} (ok);
    \draw[arr] (valid) -- node[above, font=\zihao{6}] {否} (fallback);
  \end{tikzpicture}
  \caption{问答业务流程图}
  \label{fig:qa-flow}
\end{figure}

除主成功场景外，还需要考虑三类异常。第一，检索为空。可能原因是知识库没有相关资料，或者问题表述与资料差异过大。系统应提示用户换一种问法、选择分类，或建议联系对应部门，不能强行让模型回答。第二，模型超时或返回格式错误。系统应重试一次，仍失败则返回检索到的原始片段，并记录失败日志。第三，引用校验失败。可能是模型编造了编号，也可能是证据在生成过程中被删除。系统应隐藏无效引用，降级为摘录式回答，并把这次会话标记为“需人工检查”，供管理员分析。


